\section{Design and Implementation}

\begin{frame}{Overview}
  \roadmap{0}
  \slidefig[0.66]{map}
\end{frame}

\begin{frame}{Running Example}
  \roadmap{0}
  \centering\footnotesize
  \begin{tabular}{lcc}
    \toprule
    & Example & Implementation \\
    \midrule
    Tile size & $4\times4$ & $32\times32$ \\
    Queries per batch & 4 ($q_0$--$q_3$) & 32 \\
    Inverted lists & 8 (A--H) & 512--2048 \\
    Lists per query $N_{\mathrm{probe}}$ & 2 & 1--32 \\
    Top-$k$ state & top-2 & top-32 \\
    Fine-search workers & 3 (W1--W3) & 63 \\
    Chunk limit & 2 pages & 128 pages \\
    \bottomrule
  \end{tabular}\\[2mm]
  Selections: $q_0$\{A,\,C\}, $q_1$\{A,\,B\}, $q_2$\{C,\,E\}, $q_3$\{A,\,D\}
\end{frame}

\begin{frame}{Step 1: Tiled Data Layout}
  \roadmap{1}
  \slidefig[0.3]{layout}
  \begin{itemize}\footnotesize
    \item Lists stored as \textbf{pages} of 32 vectors, dimensions padded to multiples of 32 (100 $\to$ 128)
    \item Each page has an identifier tile; padded slots carry a sentinel and are masked to $-10000$
    \item Index uploaded once to device DRAM; packed buffers addressed by page offset
  \end{itemize}
\end{frame}

\begin{frame}{Step 1: Query Batch $\times$ Page}
  \roadmap{1}
  \slidefig[0.32]{batchpage}
  \begin{itemize}\footnotesize
    \item A batch of 32 queries is the row operand, a page of 32 vectors the column operand
    \item One tile product scores every query of the batch against every vector of the page
  \end{itemize}
\end{frame}

\begin{frame}{Step 2: Coarse Search}
  \roadmap{2}
  \begin{columns}
    \column{0.62\textwidth}
    \slidefig[0.58]{coarse}
    \column{0.38\textwidth}
    \begin{itemize}\footnotesize
      \item Regular, dense: runs on the device
      \item Query tiles stay in L1, centroid blocks of 32 streamed from DRAM
      \item Running top-32 per query, so $N_{\mathrm{probe}}\le32$
      \item Example: E replaces D for $q_2$ in block 2
    \end{itemize}
  \end{columns}
\end{frame}

\begin{frame}{Step 3: Batch-Level List Union}
  \roadmap{3}
  \begin{itemize}
    \item The host collects the lists selected by all queries of a batch into one set
    \item[] \centering $\mathcal{U}_b=\bigcup_{r} \mathcal{P}_{b,r}$ \quad e.g.\ 8 selections $\to$ $\{$A,B,C,D,E$\}$
    \item Why: a page costs one tile product whether 1 or 32 queries need it
    \item Consequence: each query is also scored against the lists of the other queries
    \item With queries in dataset order, batches share few lists: each query scans \textbf{12--32$\times$} the vectors of its own lists
  \end{itemize}
\end{frame}

\begin{frame}{Step 4: Scheduling on 63 Workers}
  \roadmap{4}
  \begin{columns}
    \column{0.58\textwidth}
    \slidefig[0.55]{schedule}
    \column{0.42\textwidth}
    \begin{itemize}\footnotesize
      \item Lists longer than 128 pages split into chunks
      \item Weight = number of pages
      \item LPT: longest task to the least-loaded worker
      \item One description per worker in DRAM: (offset, page count) per task
      \item Host cost: at most 3\% of search time
    \end{itemize}
  \end{columns}
\end{frame}

\begin{frame}{Step 5: Fine Search -- Data Movement}
  \roadmap{5}
  \slidefig[0.5]{datamovement}
  \begin{itemize}\footnotesize
    \item Reader (NCRISC) streams pages over the NoC, double-buffered
    \item Compute (TRISC0--2) scores the page; writer (BRISC) stages results in DRAM
  \end{itemize}
\end{frame}

\begin{frame}{Step 5: Fine Search -- One Page}
  \roadmap{5}
  \slidefig[0.36]{page}
  \begin{itemize}\footnotesize
    \item Multiply the query tiles by the page: one $32\times32$ score tile
    \item Last page of a list: padded slots recognized by their sentinel ID and set to $-10000$
    \item Scores merged into the running top-32 of each query
  \end{itemize}
\end{frame}

\begin{frame}{Step 5: Fine Search -- Top-32 Update}
  \roadmap{5}
  \slidefig[0.5]{localtopk}
  \begin{itemize}\footnotesize
    \item Only the current page is transposed; the saved state is already kept transposed
    \item \texttt{topk\_local\_sort} orders each query column and moves scores and IDs together
  \end{itemize}
\end{frame}

\begin{frame}{Step 6: Hierarchical Top-$k$}
  \roadmap{6}
  \begin{columns}
    \column{0.6\textwidth}
    \slidefig[0.58]{staging}
    \column{0.4\textwidth}
    \begin{itemize}\footnotesize
      \item Each worker writes its partial top-32 to its own DRAM slot
      \item Completion semaphore on core $(0,0)$; ready semaphore back to each worker
      \item Core $(0,0)$ merges all slots: \textbf{final top-32 on the device}
      \item Host only reads it and keeps $k$
    \end{itemize}
  \end{columns}
\end{frame}

\begin{frame}{Step 7: Multi-Batch Execution}
  \roadmap{7}
  \slidefig[0.42]{multibatch}
  \begin{itemize}\footnotesize
    \item All batches of a workload run in one coarse and one fine invocation
    \item Workers continue with the next batch; aggregation follows behind
    \item Launch, synchronization, and host orchestration amortized over many batches
  \end{itemize}
\end{frame}
